\section{End-Term Comprehensive Examination --- Set A}

\LPUPaperHeader{End-Term Final Examination}{A}{3 Hours}{70 Marks}{CO1 to CO6 (Units 1 to 6 Comprehensive)}

\begin{center}
\sffamily\textbf{\color{AlertRed}Structure: Part A (30 Marks, 30 MCQs from Units 4, 5, 6 --- Negative Marking: $-0.25$ Marks).}\\
\textbf{\color{AlertRed}Part B (40 Marks: Attempt any 4 out of 5 Subjective Questions, 10 Marks each).}
\end{center}
\vspace{{2mm}}

\subsection*{Part A: Objective Multiple-Choice Section (30 Marks)}

\noindent\textbf{Q1.} \hfill {\small\color{LPUOrange}\textbf{[CO4 | arrays vs lists]}}\\[2pt]
Which Python library provides the \texttt{ndarray} object used for numerical arrays?

\begin{enumerate}[label=(\alph*),topsep=2pt,itemsep=1pt,partopsep=0pt,parsep=0pt]
    \item \relax Matplotlib
    \item \relax NumPy
    \item \relax SciPy
    \item \relax Pandas
\end{enumerate}
\vspace{{2mm}}

\noindent\textbf{Q2.} \hfill {\small\color{LPUOrange}\textbf{[CO4 | arrays vs lists]}}\\[2pt]
What is a key difference between a NumPy array and a Python list?

\begin{enumerate}[label=(\alph*),topsep=2pt,itemsep=1pt,partopsep=0pt,parsep=0pt]
    \item \relax A NumPy array cannot be indexed
    \item \relax A list uses less memory than an array
    \item \relax A NumPy array stores elements of a single data type
    \item \relax A list can only hold numbers
\end{enumerate}
\vspace{{2mm}}

\noindent\textbf{Q3.} \hfill {\small\color{LPUOrange}\textbf{[CO4 | arrays vs lists]}}\\[2pt]
Which function is used to create a NumPy array from a Python list?

\begin{enumerate}[label=(\alph*),topsep=2pt,itemsep=1pt,partopsep=0pt,parsep=0pt]
    \item \relax \texttt{np.list()}
    \item \relax \texttt{np.make()}
    \item \relax \texttt{np.create()}
    \item \relax \texttt{np.array()}
\end{enumerate}
\vspace{{2mm}}

\noindent\textbf{Q4.} \hfill {\small\color{LPUOrange}\textbf{[CO4 | arrays vs lists]}}\\[2pt]
Compared to Python lists, NumPy arrays are generally preferred for numerical computation because they are:

\begin{enumerate}[label=(\alph*),topsep=2pt,itemsep=1pt,partopsep=0pt,parsep=0pt]
    \item \relax Limited to one element only
    \item \relax Faster and more memory efficient
    \item \relax Unable to store integers
    \item \relax Slower but easier to read
\end{enumerate}
\vspace{{2mm}}

\noindent\textbf{Q5.} \hfill {\small\color{LPUOrange}\textbf{[CO4 | data types]}}\\[2pt]
Which attribute is used to check the data type of elements in a NumPy array?

\begin{enumerate}[label=(\alph*),topsep=2pt,itemsep=1pt,partopsep=0pt,parsep=0pt]
    \item \relax \texttt{.datatype}
    \item \relax \texttt{.dtype}
    \item \relax \texttt{.type}
    \item \relax \texttt{.kind}
\end{enumerate}
\vspace{{2mm}}

\noindent\textbf{Q6.} \hfill {\small\color{LPUOrange}\textbf{[CO4 | data types]}}\\[2pt]
What does the NumPy data type \texttt{int64} represent?

\begin{enumerate}[label=(\alph*),topsep=2pt,itemsep=1pt,partopsep=0pt,parsep=0pt]
    \item \relax A 64-character string
    \item \relax A 64-element array
    \item \relax A 64-bit floating point number
    \item \relax A 64-bit integer
\end{enumerate}
\vspace{{2mm}}

\noindent\textbf{Q7.} \hfill {\small\color{LPUOrange}\textbf{[CO4 | data types]}}\\[2pt]
Which method changes the data type of an existing NumPy array?

\begin{enumerate}[label=(\alph*),topsep=2pt,itemsep=1pt,partopsep=0pt,parsep=0pt]
    \item \relax \texttt{.astype()}
    \item \relax \texttt{.convert()}
    \item \relax \texttt{.retype()}
    \item \relax \texttt{.setdtype()}
\end{enumerate}
\vspace{{2mm}}

\noindent\textbf{Q8.} \hfill {\small\color{LPUOrange}\textbf{[CO4 | data types]}}\\[2pt]
If an array is created as \texttt{np.array([1.0, 2.0, 3.0])}, its default data type will be:

\begin{enumerate}[label=(\alph*),topsep=2pt,itemsep=1pt,partopsep=0pt,parsep=0pt]
    \item \relax \texttt{bool}
    \item \relax \texttt{str}
    \item \relax \texttt{float64}
    \item \relax \texttt{int32}
\end{enumerate}
\vspace{{2mm}}

\noindent\textbf{Q9.} \hfill {\small\color{LPUOrange}\textbf{[CO4 | array operations]}}\\[2pt]
Given \texttt{a = np.array([1, 2, 3])}, what is the result of \texttt{a + 2}?

\begin{enumerate}[label=(\alph*),topsep=2pt,itemsep=1pt,partopsep=0pt,parsep=0pt]
    \item \relax \texttt{[3, 4, 5]}
    \item \relax \texttt{[1, 2, 3, 2]}
    \item \relax \texttt{[2, 4, 6]}
    \item \relax \texttt{[3, 2, 3]}
\end{enumerate}
\vspace{{2mm}}

\noindent\textbf{Q10.} \hfill {\small\color{LPUOrange}\textbf{[CO4 | array operations]}}\\[2pt]
Which attribute returns the shape (dimensions) of a NumPy array?

\begin{enumerate}[label=(\alph*),topsep=2pt,itemsep=1pt,partopsep=0pt,parsep=0pt]
    \item \relax \texttt{.size}
    \item \relax \texttt{.shape}
    \item \relax \texttt{.dim}
    \item \relax \texttt{.length}
\end{enumerate}
\vspace{{2mm}}

\noindent\textbf{Q11.} \hfill {\small\color{LPUOrange}\textbf{[CO5 | introduction to pandas]}}\\[2pt]
What is Pandas primarily used for in Python?

\begin{enumerate}[label=(\alph*),topsep=2pt,itemsep=1pt,partopsep=0pt,parsep=0pt]
    \item \relax Managing operating system processes
    \item \relax Data manipulation and analysis
    \item \relax Compiling C++ code
    \item \relax Building 3D games
\end{enumerate}
\vspace{{2mm}}

\noindent\textbf{Q12.} \hfill {\small\color{LPUOrange}\textbf{[CO5 | introduction to pandas]}}\\[2pt]
Which is the conventional way to import the Pandas library?

\begin{enumerate}[label=(\alph*),topsep=2pt,itemsep=1pt,partopsep=0pt,parsep=0pt]
    \item \relax \texttt{include pandas as pd}
    \item \relax \texttt{import pd as pandas}
    \item \relax \texttt{import pandas as pd}
    \item \relax \texttt{import pandas.pd}
\end{enumerate}
\vspace{{2mm}}

\noindent\textbf{Q13.} \hfill {\small\color{LPUOrange}\textbf{[CO5 | introduction to pandas]}}\\[2pt]
Pandas is built on top of which numerical computing library?

\begin{enumerate}[label=(\alph*),topsep=2pt,itemsep=1pt,partopsep=0pt,parsep=0pt]
    \item \relax Matplotlib
    \item \relax TensorFlow
    \item \relax SciPy
    \item \relax NumPy
\end{enumerate}
\vspace{{2mm}}

\noindent\textbf{Q14.} \hfill {\small\color{LPUOrange}\textbf{[CO5 | series]}}\\[2pt]
A Pandas Series is best described as a:

\begin{enumerate}[label=(\alph*),topsep=2pt,itemsep=1pt,partopsep=0pt,parsep=0pt]
    \item \relax Two-dimensional labeled table
    \item \relax One-dimensional labeled array
    \item \relax Three-dimensional data cube
    \item \relax Scalar single value
\end{enumerate}
\vspace{{2mm}}

\noindent\textbf{Q15.} \hfill {\small\color{LPUOrange}\textbf{[CO5 | series]}}\\[2pt]
Which function is used to create a Pandas Series?

\begin{enumerate}[label=(\alph*),topsep=2pt,itemsep=1pt,partopsep=0pt,parsep=0pt]
    \item \relax \texttt{pd.DataFrame()}
    \item \relax \texttt{pd.array()}
    \item \relax \texttt{pd.Series()}
    \item \relax \texttt{pd.series()}
\end{enumerate}
\vspace{{2mm}}

\noindent\textbf{Q16.} \hfill {\small\color{LPUOrange}\textbf{[CO5 | series]}}\\[2pt]
What is the default index for a Series created from a list of 4 elements?

\begin{enumerate}[label=(\alph*),topsep=2pt,itemsep=1pt,partopsep=0pt,parsep=0pt]
    \item \relax Integers 1, 2, 3, 4
    \item \relax Integers 0, 1, 2, 3
    \item \relax Letters a, b, c, d
    \item \relax No index is assigned
\end{enumerate}
\vspace{{2mm}}

\noindent\textbf{Q17.} \hfill {\small\color{LPUOrange}\textbf{[CO5 | dataframe]}}\\[2pt]
A Pandas DataFrame is a:

\begin{enumerate}[label=(\alph*),topsep=2pt,itemsep=1pt,partopsep=0pt,parsep=0pt]
    \item \relax One-dimensional labeled array
    \item \relax Single scalar value
    \item \relax Plain Python dictionary
    \item \relax Two-dimensional labeled data structure
\end{enumerate}
\vspace{{2mm}}

\noindent\textbf{Q18.} \hfill {\small\color{LPUOrange}\textbf{[CO5 | dataframe]}}\\[2pt]
Which method displays the first 5 rows of a DataFrame \texttt{df} by default?

\begin{enumerate}[label=(\alph*),topsep=2pt,itemsep=1pt,partopsep=0pt,parsep=0pt]
    \item \relax \texttt{df.first()}
    \item \relax \texttt{df.head()}
    \item \relax \texttt{df.top()}
    \item \relax \texttt{df.start()}
\end{enumerate}
\vspace{{2mm}}

\noindent\textbf{Q19.} \hfill {\small\color{LPUOrange}\textbf{[CO5 | dataframe]}}\\[2pt]
Which attribute gives the number of rows and columns of a DataFrame as a tuple?

\begin{enumerate}[label=(\alph*),topsep=2pt,itemsep=1pt,partopsep=0pt,parsep=0pt]
    \item \relax \texttt{df.dimensions}
    \item \relax \texttt{df.length}
    \item \relax \texttt{df.size}
    \item \relax \texttt{df.shape}
\end{enumerate}
\vspace{{2mm}}

\noindent\textbf{Q20.} \hfill {\small\color{LPUOrange}\textbf{[CO5 | working with csv files]}}\\[2pt]
Which Pandas function reads data from a CSV file into a DataFrame?

\begin{enumerate}[label=(\alph*),topsep=2pt,itemsep=1pt,partopsep=0pt,parsep=0pt]
    \item \relax \texttt{pd.read\_csv()}
    \item \relax \texttt{pd.import\_csv()}
    \item \relax \texttt{pd.load\_csv()}
    \item \relax \texttt{pd.open\_csv()}
\end{enumerate}
\vspace{{2mm}}

\noindent\textbf{Q21.} \hfill {\small\color{LPUOrange}\textbf{[CO6 | OpenCV-based image operations]}}\\[2pt]
Which Python library is primarily used for computer vision and image processing tasks?

\begin{enumerate}[label=(\alph*),topsep=2pt,itemsep=1pt,partopsep=0pt,parsep=0pt]
    \item \relax Matplotlib
    \item \relax PyGame
    \item \relax OpenCV
    \item \relax BeautifulSoup
\end{enumerate}
\vspace{{2mm}}

\noindent\textbf{Q22.} \hfill {\small\color{LPUOrange}\textbf{[CO6 | OpenCV-based image operations]}}\\[2pt]
Which function is used to read an image from a file in OpenCV?

\begin{enumerate}[label=(\alph*),topsep=2pt,itemsep=1pt,partopsep=0pt,parsep=0pt]
    \item \relax \texttt{cv2.showImage()}
    \item \relax \texttt{cv2.loadFile()}
    \item \relax \texttt{cv2.openImage()}
    \item \relax \texttt{cv2.imread()}
\end{enumerate}
\vspace{{2mm}}

\noindent\textbf{Q23.} \hfill {\small\color{LPUOrange}\textbf{[CO6 | OpenCV-based image operations]}}\\[2pt]
In OpenCV, which color format are images loaded in by default?

\begin{enumerate}[label=(\alph*),topsep=2pt,itemsep=1pt,partopsep=0pt,parsep=0pt]
    \item \relax HSV
    \item \relax BGR
    \item \relax RGB
    \item \relax CMYK
\end{enumerate}
\vspace{{2mm}}

\noindent\textbf{Q24.} \hfill {\small\color{LPUOrange}\textbf{[CO6 | OpenCV-based image operations]}}\\[2pt]
Which OpenCV function displays an image in a window?

\begin{enumerate}[label=(\alph*),topsep=2pt,itemsep=1pt,partopsep=0pt,parsep=0pt]
    \item \relax \texttt{cv2.view()}
    \item \relax \texttt{cv2.render()}
    \item \relax \texttt{cv2.imshow()}
    \item \relax \texttt{cv2.display()}
\end{enumerate}
\vspace{{2mm}}

\noindent\textbf{Q25.} \hfill {\small\color{LPUOrange}\textbf{[CO6 | PyGame-based gaming creation]}}\\[2pt]
PyGame is a Python library mainly used for creating what type of applications?

\begin{enumerate}[label=(\alph*),topsep=2pt,itemsep=1pt,partopsep=0pt,parsep=0pt]
    \item \relax Web servers
    \item \relax Machine learning models
    \item \relax Database systems
    \item \relax Games and multimedia
\end{enumerate}
\vspace{{2mm}}

\noindent\textbf{Q26.} \hfill {\small\color{LPUOrange}\textbf{[CO6 | PyGame-based gaming creation]}}\\[2pt]
Which function must be called to start using PyGame's modules?

\begin{enumerate}[label=(\alph*),topsep=2pt,itemsep=1pt,partopsep=0pt,parsep=0pt]
    \item \relax \texttt{pygame.begin()}
    \item \relax \texttt{pygame.run()}
    \item \relax \texttt{pygame.start()}
    \item \relax \texttt{pygame.init()}
\end{enumerate}
\vspace{{2mm}}

\noindent\textbf{Q27.} \hfill {\small\color{LPUOrange}\textbf{[CO6 | PyGame-based gaming creation]}}\\[2pt]
In PyGame, which object represents the main window or drawing surface?

\begin{enumerate}[label=(\alph*),topsep=2pt,itemsep=1pt,partopsep=0pt,parsep=0pt]
    \item \relax Frame
    \item \relax Panel
    \item \relax Surface
    \item \relax Canvas
\end{enumerate}
\vspace{{2mm}}

\noindent\textbf{Q28.} \hfill {\small\color{LPUOrange}\textbf{[CO6 | PyGame-based gaming creation]}}\\[2pt]
Which PyGame concept is typically used to continuously update the game and handle events?

\begin{enumerate}[label=(\alph*),topsep=2pt,itemsep=1pt,partopsep=0pt,parsep=0pt]
    \item \relax Data pipeline
    \item \relax Game loop
    \item \relax Event stack
    \item \relax Render tree
\end{enumerate}
\vspace{{2mm}}

\noindent\textbf{Q29.} \hfill {\small\color{LPUOrange}\textbf{[CO6 | web-scraping-based url access]}}\\[2pt]
Which Python library is commonly used to send HTTP requests to access a URL?

\begin{enumerate}[label=(\alph*),topsep=2pt,itemsep=1pt,partopsep=0pt,parsep=0pt]
    \item \relax \texttt{requests}
    \item \relax \texttt{turtle}
    \item \relax \texttt{numpy}
    \item \relax \texttt{pygame}
\end{enumerate}
\vspace{{2mm}}

\noindent\textbf{Q30.} \hfill {\small\color{LPUOrange}\textbf{[CO6 | web-scraping-based url access]}}\\[2pt]
Which library is widely used to parse HTML content during web scraping?

\begin{enumerate}[label=(\alph*),topsep=2pt,itemsep=1pt,partopsep=0pt,parsep=0pt]
    \item \relax Seaborn
    \item \relax BeautifulSoup
    \item \relax PyGame
    \item \relax OpenCV
\end{enumerate}
\vspace{{2mm}}

\newpage
\subsection*{Part B: Descriptive Analytical \& Programming Section (40 Marks)}

\sffamily\textbf{Note: Attempt any FOUR questions from this section. Each question carries 10 marks.}
\vspace{{2mm}}

\questionitem{31}{Design a complete Python class \texttt{MatrixOperations} utilizing NumPy that performs matrix inversion, determinant calculation via LU decomposition, and solves a system of linear equations $Ax = b$ with error handling for singular matrices.}{10}{CO4}{L4 Analyze}
\vspace{3mm}

\questionitem{32}{Develop an end-to-end Pandas and Matplotlib data analytics pipeline to ingest a web transaction log, extract hourly traffic volume, compute 7-day moving averages, and plot a multi-curve trend with confidence intervals.}{10}{CO5}{L3 Apply}
\vspace{3mm}

\questionitem{33}{Design a real-time computer vision algorithm using OpenCV to track a colored object in a video stream, compute its bounding box centroid coordinates, and draw trajectory overlays.}{10}{CO6}{L3 Apply}
\vspace{3mm}

\questionitem{34}{Explain the complete Pygame event pump architecture. Construct an interactive game loop with a movable player sprite, wall collision boundaries, and frame-rate regulation at 60 FPS.}{10}{CO6}{L3 Apply}
\vspace{3mm}

\questionitem{35}{Design a resilient multi-threaded web scraper using \texttt{concurrent.futures}, \texttt{requests}, and \texttt{BeautifulSoup} to scrape research paper abstracts from an archive portal, handle HTTP timeouts, and persist structured JSON records.}{10}{CO6}{L4 Analyze}
\vspace{3mm}

\newpage
\subsection*{Official Solutions and Marking Rubrics --- End-Term Set A}

\subsubsection*{Part A Answer Key}

\begin{center}
\begin{tabularx}{0.96\textwidth}{c c X c c X}
\toprule
\textbf{Q\#} & \textbf{Ans} & \textbf{Core Topic} & \textbf{Q\#} & \textbf{Ans} & \textbf{Core Topic} \\
\midrule
1 & \textbf{B} & ndarray contiguous C-order buffer & 16 & \textbf{C} & Groupby split-apply-combine \\
2 & \textbf{A} & Vectorized broadcasting compatibility & 17 & \textbf{A} & Resample timeseries frequency \\
3 & \textbf{D} & Stride calculation for view & 18 & \textbf{D} & Merge join types inner/outer \\
4 & \textbf{C} & Fancy indexing produces copy & 19 & \textbf{B} & Matplotlib subplots OOP model \\
5 & \textbf{A} & Universal function ufunc & 20 & \textbf{C} & Histogram density normalization \\
6 & \textbf{B} & Reduction axis collapses dimension & 21 & \textbf{A} & OpenCV loads in BGR order \\
7 & \textbf{C} & Reshape total elements invariant & 22 & \textbf{B} & Color conversion COLOR\_BGR2GRAY \\
8 & \textbf{D} & Flatten copy vs ravel view & 23 & \textbf{D} & Gaussian blur 2D convolution \\
9 & \textbf{B} & Matrix product @ vs * Hadamard & 24 & \textbf{C} & Canny double threshold hysteresis \\
10 & \textbf{A} & Broadcasting dimensions alignment & 25 & \textbf{A} & Contour retrieval external \\
11 & \textbf{C} & Series 1D labeled array & 26 & \textbf{D} & Pygame canonical event loop \\
12 & \textbf{B} & DataFrame 2D tabular data & 27 & \textbf{B} & Clock tick locks FPS rate \\
13 & \textbf{A} & .loc inclusive vs .iloc exclusive & 28 & \textbf{C} & Requests HTTP status checking \\
14 & \textbf{D} & SettingWithCopyWarning prevention & 29 & \textbf{A} & BeautifulSoup DOM parser \\
15 & \textbf{B} & Missing value detection isna & 30 & \textbf{D} & Scikit-learn fit on train set only \\
\bottomrule
\end{tabularx}
\end{center}
\vspace{{4mm}}

\subsubsection*{Part B Comprehensive Solutions \& Marking Schemes}

\begin{solutionbox}[Question 31 Solution]
\begin{pycode}
import numpy as np\\
~\\
class MatrixOperations:\\
\_\_ESCTOKEN\_7\_\_ def \_\_ESCTOKEN\_0\_\_\_\_ESCTOKEN\_0\_\_init\_\_ESCTOKEN\_0\_\_\_\_ESCTOKEN\_0\_\_(self, A):\\
\_\_ESCTOKEN\_7\_\_ \_\_ESCTOKEN\_7\_\_ self.A = np.array(A, dtype=float)\\
\_\_ESCTOKEN\_7\_\_ \_\_ESCTOKEN\_7\_\_ if self.A.ndim != 2 or self.A.shape[0] != self.A.shape[1]:\\
\_\_ESCTOKEN\_7\_\_ \_\_ESCTOKEN\_7\_\_ \_\_ESCTOKEN\_7\_\_ raise ValueError("Matrix must be 2D square")\\
~\\
\_\_ESCTOKEN\_7\_\_ def determinant(self):\\
\_\_ESCTOKEN\_7\_\_ \_\_ESCTOKEN\_7\_\_ return np.linalg.det(self.A)\\
~\\
\_\_ESCTOKEN\_7\_\_ def solve\_\_ESCTOKEN\_0\_\_system(self, b):\\
\_\_ESCTOKEN\_7\_\_ \_\_ESCTOKEN\_7\_\_ b = np.array(b, dtype=float)\\
\_\_ESCTOKEN\_7\_\_ \_\_ESCTOKEN\_7\_\_ det = self.determinant()\\
\_\_ESCTOKEN\_7\_\_ \_\_ESCTOKEN\_7\_\_ if np.isclose(det, 0.0):\\
\_\_ESCTOKEN\_7\_\_ \_\_ESCTOKEN\_7\_\_ \_\_ESCTOKEN\_7\_\_ raise np.linalg.LinAlgError("Matrix is singular (det=0); no unique solution exists")\\
\_\_ESCTOKEN\_7\_\_ \_\_ESCTOKEN\_7\_\_ return np.linalg.solve(self.A, b)
\end{pycode}
\textbf{Complexity Analysis}: LU factorization requires $\mathcal{O}(n^3)$ floating-point operations. The \texttt{np.isclose} check prevents numerical instability artifacts in near-singular matrices.
\begin{markingrubric}
$\bullet$ Square matrix validation and OOP design: \textbf{3 Marks}\\$\bullet$ Determinant and LinAlgError exception handling: \textbf{4 Marks}\\$\bullet$ Asymptotic complexity $\mathcal{O}(n^3)$ derivation: \textbf{3 Marks}
\end{markingrubric}
\end{solutionbox}

\begin{solutionbox}[Question 32 Solution]
\begin{pycode}
import pandas as pd\\
import matplotlib.pyplot as plt\\
~\\
def analyze\_\_ESCTOKEN\_0\_\_traffic(csv\_\_ESCTOKEN\_0\_\_path):\\
\_\_ESCTOKEN\_7\_\_ df = pd.read\_\_ESCTOKEN\_0\_\_csv(csv\_\_ESCTOKEN\_0\_\_path, parse\_\_ESCTOKEN\_0\_\_dates=['timestamp'])\\
\_\_ESCTOKEN\_7\_\_ df.set\_\_ESCTOKEN\_0\_\_index('timestamp', inplace=True)\\
\_\_ESCTOKEN\_7\_\_ hourly = df.resample('H').size().rename('requests')\\
\_\_ESCTOKEN\_7\_\_ rolling\_\_ESCTOKEN\_0\_\_7d = hourly.rolling(24 * 7).mean()\\
\_\_ESCTOKEN\_7\_\_ std\_\_ESCTOKEN\_0\_\_7d = hourly.rolling(24 * 7).std()\\
~\\
\_\_ESCTOKEN\_7\_\_ plt.figure(figsize=(14, 6))\\
\_\_ESCTOKEN\_7\_\_ plt.plot(hourly.index, hourly, alpha=0.3, label='Hourly Traffic')\\
\_\_ESCTOKEN\_7\_\_ plt.plot(rolling\_\_ESCTOKEN\_0\_\_7d.index, rolling\_\_ESCTOKEN\_0\_\_7d, color='navy', label='7D Moving Average')\\
\_\_ESCTOKEN\_7\_\_ plt.fill\_\_ESCTOKEN\_0\_\_between(rolling\_\_ESCTOKEN\_0\_\_7d.index, rolling\_\_ESCTOKEN\_0\_\_7d - 2*std\_\_ESCTOKEN\_0\_\_7d, rolling\_\_ESCTOKEN\_0\_\_7d + 2*std\_\_ESCTOKEN\_0\_\_7d,\\
\_\_ESCTOKEN\_7\_\_ \_\_ESCTOKEN\_7\_\_ \_\_ESCTOKEN\_7\_\_ \_\_ESCTOKEN\_7\_\_ \_\_ESCTOKEN\_7\_\_ color='blue', alpha=0.15, label='95\_\_ESCTOKEN\_2\_\_ Confidence Band')\\
\_\_ESCTOKEN\_7\_\_ plt.title("Web Traffic Trends")\\
\_\_ESCTOKEN\_7\_\_ plt.legend(); plt.tight\_\_ESCTOKEN\_0\_\_layout()\\
\_\_ESCTOKEN\_7\_\_ plt.savefig("traffic.png")
\end{pycode}
\begin{markingrubric}
$\bullet$ Ingestion, timestamp parsing, and resampling: \textbf{3 Marks}\\$\bullet$ Rolling statistics computation: \textbf{3 Marks}\\$\bullet$ Multi-curve Matplotlib plot with shaded confidence region: \textbf{4 Marks}
\end{markingrubric}
\end{solutionbox}

\begin{solutionbox}[Question 33 Solution]
\begin{pycode}
import cv2\\
import numpy as np\\
~\\
def track\_\_ESCTOKEN\_0\_\_colored\_\_ESCTOKEN\_0\_\_object(frame):\\
\_\_ESCTOKEN\_7\_\_ hsv = cv2.cvtColor(frame, cv2.COLOR\_\_ESCTOKEN\_0\_\_BGR2HSV)\\
\_\_ESCTOKEN\_7\_\_ lower\_\_ESCTOKEN\_0\_\_blue = np.array([100, 150, 50])\\
\_\_ESCTOKEN\_7\_\_ upper\_\_ESCTOKEN\_0\_\_blue = np.array([140, 255, 255])\\
\_\_ESCTOKEN\_7\_\_ mask = cv2.inRange(hsv, lower\_\_ESCTOKEN\_0\_\_blue, upper\_\_ESCTOKEN\_0\_\_blue)\\
\_\_ESCTOKEN\_7\_\_ mask = cv2.erode(mask, None, iterations=2)\\
\_\_ESCTOKEN\_7\_\_ mask = cv2.dilate(mask, None, iterations=2)\\
~\\
\_\_ESCTOKEN\_7\_\_ contours, \_\_ESCTOKEN\_0\_\_ = cv2.findContours(mask, cv2.RETR\_\_ESCTOKEN\_0\_\_EXTERNAL, cv2.CHAIN\_\_ESCTOKEN\_0\_\_APPROX\_\_ESCTOKEN\_0\_\_SIMPLE)\\
\_\_ESCTOKEN\_7\_\_ if contours:\\
\_\_ESCTOKEN\_7\_\_ \_\_ESCTOKEN\_7\_\_ c = max(contours, key=cv2.contourArea)\\
\_\_ESCTOKEN\_7\_\_ \_\_ESCTOKEN\_7\_\_ ((x, y), radius) = cv2.minEnclosingCircle(c)\\
\_\_ESCTOKEN\_7\_\_ \_\_ESCTOKEN\_7\_\_ M = cv2.moments(c)\\
\_\_ESCTOKEN\_7\_\_ \_\_ESCTOKEN\_7\_\_ if M["m00"] > 0:\\
\_\_ESCTOKEN\_7\_\_ \_\_ESCTOKEN\_7\_\_ \_\_ESCTOKEN\_7\_\_ cx = int(M["m10"] / M["m00"])\\
\_\_ESCTOKEN\_7\_\_ \_\_ESCTOKEN\_7\_\_ \_\_ESCTOKEN\_7\_\_ cy = int(M["m01"] / M["m00"])\\
\_\_ESCTOKEN\_7\_\_ \_\_ESCTOKEN\_7\_\_ \_\_ESCTOKEN\_7\_\_ cv2.circle(frame, (cx, cy), 5, (0, 0, 255), -1)\\
\_\_ESCTOKEN\_7\_\_ return frame
\end{pycode}
\begin{markingrubric}
$\bullet$ Color space conversion to HSV and thresholding: \textbf{3 Marks}\\$\bullet$ Morphological filtering (erosion/dilation): \textbf{3 Marks}\\$\bullet$ Centroid calculation via spatial moments ($m_{10}/m_{00}$): \textbf{4 Marks}
\end{markingrubric}
\end{solutionbox}

\begin{solutionbox}[Question 34 Solution]
\begin{pycode}
import pygame, sys\\
~\\
pygame.init()\\
screen = pygame.display.set\_\_ESCTOKEN\_0\_\_mode((800, 600))\\
clock = pygame.time.Clock()\\
player = pygame.Rect(380, 280, 40, 40)\\
speed = 5\\
~\\
while True:\\
\_\_ESCTOKEN\_7\_\_ for event in pygame.event.get():\\
\_\_ESCTOKEN\_7\_\_ \_\_ESCTOKEN\_7\_\_ if event.type == pygame.QUIT:\\
\_\_ESCTOKEN\_7\_\_ \_\_ESCTOKEN\_7\_\_ \_\_ESCTOKEN\_7\_\_ pygame.quit(); sys.exit()\\
~\\
\_\_ESCTOKEN\_7\_\_ keys = pygame.key.get\_\_ESCTOKEN\_0\_\_pressed()\\
\_\_ESCTOKEN\_7\_\_ if keys[pygame.K\_\_ESCTOKEN\_0\_\_LEFT]:  player.x -= speed\\
\_\_ESCTOKEN\_7\_\_ if keys[pygame.K\_\_ESCTOKEN\_0\_\_RIGHT]: player.x += speed\\
\_\_ESCTOKEN\_7\_\_ if keys[pygame.K\_\_ESCTOKEN\_0\_\_UP]:    player.y -= speed\\
\_\_ESCTOKEN\_7\_\_ if keys[pygame.K\_\_ESCTOKEN\_0\_\_DOWN]:  player.y += speed\\
~\\
\_\_ESCTOKEN\_7\_\_ \_\_ESCTOKEN\_1\_\_ Boundary clipping:\\
\_\_ESCTOKEN\_7\_\_ player.clamp\_\_ESCTOKEN\_0\_\_ip(screen.get\_\_ESCTOKEN\_0\_\_rect())\\
~\\
\_\_ESCTOKEN\_7\_\_ screen.fill((30, 30, 30))\\
\_\_ESCTOKEN\_7\_\_ pygame.draw.rect(screen, (0, 200, 255), player)\\
\_\_ESCTOKEN\_7\_\_ pygame.display.flip()\\
\_\_ESCTOKEN\_7\_\_ clock.tick(60)
\end{pycode}
\begin{markingrubric}
$\bullet$ Pygame initialization and surface management: \textbf{2 Marks}\\$\bullet$ Event pump handling and graceful shutdown: \textbf{3 Marks}\\$\bullet$ Boundary collision clamp and 60 FPS clock regulator: \textbf{5 Marks}
\end{markingrubric}
\end{solutionbox}

\begin{solutionbox}[Question 35 Solution]
\begin{pycode}
import requests, json\\
from bs4 import BeautifulSoup\\
from concurrent.futures import ThreadPoolExecutor\\
~\\
def fetch\_\_ESCTOKEN\_0\_\_abstract(paper\_\_ESCTOKEN\_0\_\_id):\\
\_\_ESCTOKEN\_7\_\_ url = f"https://example-archive.org/paper/\_\_ESCTOKEN\_5\_\_paper\_\_ESCTOKEN\_0\_\_id\_\_ESCTOKEN\_6\_\_"\\
\_\_ESCTOKEN\_7\_\_ try:\\
\_\_ESCTOKEN\_7\_\_ \_\_ESCTOKEN\_7\_\_ resp = requests.get(url, timeout=5, headers=\_\_ESCTOKEN\_5\_\_'User-Agent': 'Bot/1.0'\_\_ESCTOKEN\_6\_\_)\\
\_\_ESCTOKEN\_7\_\_ \_\_ESCTOKEN\_7\_\_ if resp.status\_\_ESCTOKEN\_0\_\_code == 200:\\
\_\_ESCTOKEN\_7\_\_ \_\_ESCTOKEN\_7\_\_ \_\_ESCTOKEN\_7\_\_ soup = BeautifulSoup(resp.text, 'html.parser')\\
\_\_ESCTOKEN\_7\_\_ \_\_ESCTOKEN\_7\_\_ \_\_ESCTOKEN\_7\_\_ title = soup.find('h1', class\_\_ESCTOKEN\_0\_\_='title').get\_\_ESCTOKEN\_0\_\_text(strip=True)\\
\_\_ESCTOKEN\_7\_\_ \_\_ESCTOKEN\_7\_\_ \_\_ESCTOKEN\_7\_\_ abstract = soup.find('div', class\_\_ESCTOKEN\_0\_\_='abstract').get\_\_ESCTOKEN\_0\_\_text(strip=True)\\
\_\_ESCTOKEN\_7\_\_ \_\_ESCTOKEN\_7\_\_ \_\_ESCTOKEN\_7\_\_ return \_\_ESCTOKEN\_5\_\_'id': paper\_\_ESCTOKEN\_0\_\_id, 'title': title, 'abstract': abstract\_\_ESCTOKEN\_6\_\_\\
\_\_ESCTOKEN\_7\_\_ except requests.exceptions.RequestException:\\
\_\_ESCTOKEN\_7\_\_ \_\_ESCTOKEN\_7\_\_ return None\\
~\\
def run\_\_ESCTOKEN\_0\_\_scraper(paper\_\_ESCTOKEN\_0\_\_ids):\\
\_\_ESCTOKEN\_7\_\_ with ThreadPoolExecutor(max\_\_ESCTOKEN\_0\_\_workers=5) as executor:\\
\_\_ESCTOKEN\_7\_\_ \_\_ESCTOKEN\_7\_\_ results = list(executor.map(fetch\_\_ESCTOKEN\_0\_\_abstract, paper\_\_ESCTOKEN\_0\_\_ids))\\
\_\_ESCTOKEN\_7\_\_ valid = [r for r in results if r is not None]\\
\_\_ESCTOKEN\_7\_\_ with open('papers.json', 'w') as f:\\
\_\_ESCTOKEN\_7\_\_ \_\_ESCTOKEN\_7\_\_ json.dump(valid, f, indent=2)
\end{pycode}
\begin{markingrubric}
$\bullet$ Resilient HTTP requests with timeouts and exception catching: \textbf{3 Marks}\\$\bullet$ Concurrent execution via ThreadPoolExecutor: \textbf{4 Marks}\\$\bullet$ DOM parsing and structured JSON persistence: \textbf{3 Marks}
\end{markingrubric}
\end{solutionbox}
